\chapter{Theoretische Grundlagen}\label{chapter:grundlagen}

Dieses Kapitel stellt die für die Arbeit benötigten Grundlagen zusammen.
Ausgehend von den Stoffeigenschaften des Wasserstoffs werden zunächst die
Grundbegriffe der vorgemischten Verbrennung eingeführt
(Abschnitt \ref{sec:vormischverbrennung}). Es folgen die Reaktionskinetik des
Systems \ce{H2}/\ce{O2} (Abschnitt \ref{sec:kinetik}), die laminare
Flammengeschwindigkeit als zentrale Kenngröße (Abschnitt \ref{sec:laminar})
sowie der Einfluss von Flammenstreckung und thermodiffusiver Instabilität
(Abschnitt \ref{sec:streckung}). Abschnitt \ref{sec:turbulent} behandelt die
turbulente Vormischverbrennung, Abschnitt \ref{sec:flashback} die Mechanismen
des Flammenrückschlags. Abschnitt \ref{sec:cantera} beschreibt schließlich das
numerische Modell.

\section{Wasserstoff als Brennstoff}\label{sec:wasserstoff}

Wasserstoff unterscheidet sich in nahezu allen verbrennungsrelevanten
Eigenschaften deutlich von Methan, dem Hauptbestandteil von Erdgas.
Tabelle \ref{tab:stoffwerte} stellt die wichtigsten Kenngrößen gegenüber.

\begin{table}[htbp]
\begin{center}
\caption[Verbrennungstechnische Eigenschaften von Wasserstoff und Methan]{Verbrennungstechnische Eigenschaften von Wasserstoff und Methan bei einem Druck von 1\,bar und einer Temperatur von 298\,K. Zusammengestellt nach \cite{Law_2006, Warnatz_2001, Verhelst_2009}.}
\label{tab:stoffwerte}
\begin{tabular}{lccc}
\toprule[1.2pt]
\textbf{Eigenschaft} & \textbf{Einheit} & \textbf{\ce{H2}} & \textbf{\ce{CH4}} \\
\midrule[1.2pt]\midrule[1.0pt]
Molare Masse                        & g/mol             & 2,016  & 16,04 \\
Unterer Heizwert (massebezogen)     & MJ/kg             & 120    & 50,0  \\
Unterer Heizwert (volumenbezogen)   & MJ/m$^3$          & 10,8   & 35,9  \\
Stöchiometrischer Luftbedarf        & kg/kg             & 34,3   & 17,2  \\
Zündgrenzen in Luft                 & Vol.-\%           & 4--75  & 5--15 \\
Mindestzündenergie                  & mJ                & 0,017  & 0,28  \\
Löschabstand                        & mm                & 0,64   & 2,03  \\
Diffusionskoeffizient in Luft       & cm$^2$/s          & 0,61   & 0,16  \\
Max.\ lam.\ Flammengeschwindigkeit  & m/s               & ca.\ 2,9 & ca.\ 0,40 \\
zugehöriges $\varphi$               & --                & ca.\ 1,7 & ca.\ 1,1 \\
Adiabate Flammentemperatur ($\varphi=1$) & K            & ca.\ 2390 & ca.\ 2225 \\
\bottomrule
\end{tabular}
\end{center}
\end{table}

Für die Frage des Flammenrückschlags sind vier dieser Eigenschaften von
besonderer Bedeutung. Erstens liegt die maximale laminare
Flammengeschwindigkeit rund siebenmal höher als bei Methan, sodass die Flamme
bereits bei deutlich höheren Strömungsgeschwindigkeiten stromauf laufen kann.
Zweitens erlaubt der weite Zündbereich einen Betrieb bei sehr kleinen
Äquivalenzverhältnissen, in denen Kohlenwasserstoffflammen bereits verlöschen
würden; magere Wasserstoffflammen sind gerade dort jedoch besonders anfällig
für zellulare Instabilitäten. Drittens bedeutet der geringe Löschabstand, dass
die Flamme auch in engen wandnahen Bereichen und in Spalten überleben kann, in
denen eine Methanflamme durch Wärmeverlust an die Wand erlischt. Viertens führt
der hohe Diffusionskoeffizient zu einer Lewis-Zahl deutlich unterhalb von eins,
was die in Abschnitt \ref{sec:streckung} beschriebenen thermodiffusiven
Effekte hervorruft.

Der volumenbezogene Heizwert von Wasserstoff beträgt nur etwa
30\,\% desjenigen von Methan. Bei gleicher thermischer Leistung muss folglich
ein deutlich größerer Brennstoffvolumenstrom bereitgestellt werden, was
unmittelbare Konsequenzen für die Auslegung der Zuleitungen und der
Mischungszone hat.

\section{Grundbegriffe der vorgemischten Verbrennung}\label{sec:vormischverbrennung}

\subsection{Äquivalenzverhältnis}\label{subsec:phi}

Die vollständige Oxidation von Wasserstoff mit Luft lässt sich unter der
üblichen Näherung, dass trockene Luft aus 21\,Vol.-\% Sauerstoff und
79\,Vol.-\% Stickstoff besteht, durch die Bruttoreaktionsgleichung

\begin{equation}
	\label{eq:bruttoreaktion}
	\ce{H2} + \frac{1}{2}\left(\ce{O2} + 3{,}76\,\ce{N2}\right)
	\longrightarrow \ce{H2O} + 1{,}88\,\ce{N2}
\end{equation}

beschreiben. Das Äquivalenzverhältnis $\varphi$ setzt das tatsächlich
vorliegende Brennstoff-Oxidator-Verhältnis ins Verhältnis zum
stöchiometrischen Wert,

\begin{equation}
	\label{eq:phi}
	\varphi = \frac{\left(\dot{m}_\mathrm{B}/\dot{m}_\mathrm{L}\right)}
	               {\left(\dot{m}_\mathrm{B}/\dot{m}_\mathrm{L}\right)_\mathrm{st}}
	        = \frac{1}{\lambda},
\end{equation}

wobei $\dot{m}_\mathrm{B}$ und $\dot{m}_\mathrm{L}$ die Massenströme von
Brennstoff und Luft bezeichnen und $\lambda$ die in der Anlagentechnik
gebräuchliche Luftzahl ist. Für $\varphi < 1$ liegt ein mageres, für
$\varphi > 1$ ein fettes Gemisch vor. Im stöchiometrischen Fall
($\varphi = 1$) beträgt der Wasserstoffanteil im Gemisch 29,6\,Vol.-\%.

Da an dem in dieser Arbeit betrachteten Prüfstand die Volumenströme von
Wasserstoff und Luft über Massendurchflussregler eingestellt werden, ist es
zweckmäßig, $\varphi$ direkt aus den Normvolumenströmen zu bestimmen:

\begin{equation}
	\label{eq:phi_volumen}
	\varphi = \frac{\dot{V}_{\ce{H2}}}{\dot{V}_\mathrm{Luft}}
	          \cdot \frac{1}{\left(\dot{V}_{\ce{H2}}/\dot{V}_\mathrm{Luft}\right)_\mathrm{st}}
	        = \frac{\dot{V}_{\ce{H2}}}{\dot{V}_\mathrm{Luft}} \cdot 2{,}38.
\end{equation}

\subsection{Adiabate Flammentemperatur}\label{subsec:adiabat}

Die adiabate Flammentemperatur $T_\mathrm{ad}$ ergibt sich aus der Bedingung,
dass die Enthalpie des Systems bei isobarer, adiabater Umsetzung erhalten
bleibt:

\begin{equation}
	\label{eq:adiabat}
	h\!\left(T_\mathrm{u}, Y_{k,\mathrm{u}}\right)
	= h\!\left(T_\mathrm{ad}, Y_{k,\mathrm{b}}\right).
\end{equation}

Hierin bezeichnen die Indizes $\mathrm{u}$ und $\mathrm{b}$ das unverbrannte
beziehungsweise verbrannte Gas und $Y_k$ die Massenbrüche der beteiligten
Spezies. Bei hohen Temperaturen ist die Dissoziation der Produkte zu
berücksichtigen, sodass $T_\mathrm{ad}$ nicht aus einer einfachen
Enthalpiebilanz mit vollständiger Umsetzung folgt, sondern über die Bedingung
des chemischen Gleichgewichts bestimmt werden muss \cite{Warnatz_2001}. Das
Maximum von $T_\mathrm{ad}$ liegt für Wasserstoff-Luft-Gemische geringfügig auf
der fetten Seite von $\varphi = 1$.

\section{Chemische Reaktionskinetik}\label{sec:kinetik}

\subsection{Elementarreaktionen und Reaktionsgeschwindigkeit}\label{subsec:elementarreaktionen}

Die Bruttoreaktionsgleichung \eqref{eq:bruttoreaktion} gibt lediglich die
Bilanz zwischen Edukten und Produkten wieder. Der tatsächliche Ablauf der
Oxidation setzt sich aus einer Vielzahl von Elementarreaktionen zwischen
Radikalen und stabilen Spezies zusammen. Für eine Elementarreaktion $i$ der
Form

\begin{equation}
	\label{eq:elementarreaktion}
	\sum_{k=1}^{K} \nu_{ki}^{\prime}\,\mathcal{M}_k
	\;\rightleftharpoons\;
	\sum_{k=1}^{K} \nu_{ki}^{\prime\prime}\,\mathcal{M}_k
\end{equation}

folgt die Netto-Reaktionsrate aus dem Massenwirkungsgesetz zu

\begin{equation}
	\label{eq:reaktionsrate}
	q_i = k_{\mathrm{f},i} \prod_{k=1}^{K} \left[X_k\right]^{\nu_{ki}^{\prime}}
	    - k_{\mathrm{r},i} \prod_{k=1}^{K} \left[X_k\right]^{\nu_{ki}^{\prime\prime}},
\end{equation}

wobei $\left[X_k\right]$ die molare Konzentration der Spezies $k$ bezeichnet.
Die Geschwindigkeitskoeffizienten $k_{\mathrm{f},i}$ werden üblicherweise in
der erweiterten Arrhenius-Form

\begin{equation}
	\label{eq:arrhenius}
	k_{\mathrm{f},i}(T) = A_i\,T^{b_i}\,\exp\!\left(-\frac{E_{\mathrm{a},i}}{R_\mathrm{m} T}\right)
\end{equation}

angegeben. Die Bildungsrate der Spezies $k$ ergibt sich durch Summation über
alle $I$ Reaktionen:

\begin{equation}
	\label{eq:speziesquellterm}
	\dot{\omega}_k = \sum_{i=1}^{I}
	\left(\nu_{ki}^{\prime\prime} - \nu_{ki}^{\prime}\right) q_i.
\end{equation}

\subsection{Reaktionspfade der Wasserstoffoxidation}\label{subsec:reaktionspfade}

Das Reaktionssystem \ce{H2}/\ce{O2} umfasst je nach Mechanismus etwa
acht bis zehn Spezies (\ce{H2}, \ce{O2}, \ce{H2O}, \ce{H}, \ce{O}, \ce{OH},
\ce{HO2}, \ce{H2O2} sowie inerte Stoßpartner) und zwanzig bis dreißig
Elementarreaktionen. Es ist damit das am besten untersuchte Verbrennungssystem
überhaupt und bildet zugleich das Grundgerüst jedes
Kohlenwasserstoffmechanismus \cite{Law_2006}.

Von zentraler Bedeutung ist die Kettenverzweigungsreaktion

\begin{equation}
	\label{eq:kettenverzweigung}
	\ce{H} + \ce{O2} \rightleftharpoons \ce{O} + \ce{OH},
\end{equation}

bei der aus einem Radikal zwei entstehen. Sie besitzt eine vergleichsweise hohe
Aktivierungsenergie und bestimmt deshalb maßgeblich die Temperaturabhängigkeit
der Reaktionsgeschwindigkeit. In Konkurrenz dazu steht die
Dreistoßrekombination

\begin{equation}
	\label{eq:kettenabbruch}
	\ce{H} + \ce{O2} + \ce{M} \rightleftharpoons \ce{HO2} + \ce{M},
\end{equation}

die als Kettenabbruch wirkt, da das entstehende \ce{HO2} unter mageren
Bedingungen vergleichsweise reaktionsträge ist. Da Reaktion
\eqref{eq:kettenabbruch} als Dreistoßreaktion druckabhängig ist, Reaktion
\eqref{eq:kettenverzweigung} dagegen nicht, verschiebt sich das Verhältnis
beider Pfade mit steigendem Druck zugunsten des Kettenabbruchs. Aus der
Gleichgewichtsbedingung beider Raten folgt die zweite Explosionsgrenze des
Systems \ce{H2}/\ce{O2}. Dieser Zusammenhang erklärt unmittelbar, weshalb die
laminare Flammengeschwindigkeit von Wasserstoff mit steigendem Druck abnimmt
\cite{Law_2006, Burke_2012}.

Die weitere Umsetzung erfolgt über die Kettenfortpflanzungsreaktionen

\begin{align}
	\label{eq:kettenfortpflanzung}
	\ce{O} + \ce{H2}  &\rightleftharpoons \ce{H} + \ce{OH}, \\
	\ce{OH} + \ce{H2} &\rightleftharpoons \ce{H} + \ce{H2O},
\end{align}

wobei die zweite Reaktion den Hauptanteil der Wasserbildung trägt.

\subsection{Verfügbare Reaktionsmechanismen}\label{subsec:mechanismen}

Für die vorliegende Arbeit kommen mehrere validierte Mechanismen in Betracht.
Der Mechanismus von Ó Conaire et al.\ \cite{OConaire_2004} ist über einen
weiten Bereich von Druck, Temperatur und Äquivalenzverhältnis validiert und
wird häufig als Referenz herangezogen. Li et al.\ \cite{Li_2004} und
Burke et al.\ \cite{Burke_2012} legen den Schwerpunkt auf erhöhte Drücke,
wobei Burke et al.\ insbesondere die Druckabhängigkeit der Reaktion
\eqref{eq:kettenabbruch} überarbeitet haben. Kéromnès et al.\
\cite{Keromnes_2013} erweitern die Validierungsbasis auf Synthesegasgemische.
Der San-Diego-Mechanismus \cite{SanDiego_Mech} zeichnet sich durch einen
bewusst reduzierten Umfang bei gleichzeitig guter Vorhersagequalität aus.

Da die Unterschiede zwischen diesen Mechanismen in der berechneten laminaren
Flammengeschwindigkeit je nach Betriebspunkt mehrere Prozent betragen können,
wird in dieser Arbeit ein Vergleich mehrerer Mechanismen durchgeführt, um die
resultierende Modellunsicherheit zu quantifizieren.

\section{Die laminare Vormischflamme}\label{sec:laminar}

\subsection{Flammenstruktur und laminare Flammengeschwindigkeit}\label{subsec:sl}

Eine frei propagierende, adiabate und eindimensionale Vormischflamme breitet
sich relativ zum unverbrannten Gas mit der laminaren Flammengeschwindigkeit
$s_\mathrm{L}$ aus. Sie ist eine reine Stoffeigenschaft des Gemisches und hängt
ausschließlich von Zusammensetzung, Vorwärmtemperatur $T_\mathrm{u}$ und Druck
$p$ ab.

Die Flamme lässt sich in eine Vorwärmzone, in der das Frischgas durch
Wärmeleitung aus der Reaktionszone aufgeheizt wird, und eine dünne
Reaktionszone unterteilen, in der die Wärmefreisetzung stattfindet. Aus der
Betrachtung des Gleichgewichts zwischen diffusivem Wärmetransport und
chemischer Umsatzrate folgt die auf Mallard und Le Chatelier zurückgehende
Skalierung

\begin{equation}
	\label{eq:mallard}
	s_\mathrm{L} \sim \sqrt{a \cdot \dot{\omega}},
	\qquad
	a = \frac{\lambda}{\rho\,c_p},
\end{equation}

mit der Temperaturleitfähigkeit $a$ und einer charakteristischen
Reaktionsrate $\dot{\omega}$. Die zugehörige laminare Flammendicke wird in
dieser Arbeit über den maximalen Temperaturgradienten definiert,

\begin{equation}
	\label{eq:flammendicke}
	\delta_\mathrm{L} = \frac{T_\mathrm{b} - T_\mathrm{u}}
	{\left|\mathrm{d}T/\mathrm{d}x\right|_\mathrm{max}},
\end{equation}

da diese Definition unmittelbar aus dem berechneten Temperaturprofil folgt und
im Gegensatz zur diffusiven Abschätzung $\delta_\mathrm{D} = a/s_\mathrm{L}$
keine Annahme über konstante Stoffwerte erfordert \cite{Poinsot_2011}.

\subsection{Einflussgrößen}\label{subsec:einflussgroessen}

Die Abhängigkeit der laminaren Flammengeschwindigkeit vom
Äquivalenzverhältnis weist ein ausgeprägtes Maximum auf. Bei
Kohlenwasserstoffen liegt dieses nahe der Stöchiometrie
($\varphi \approx 1{,}05\ldots1{,}1$), da dort die Flammentemperatur maximal
wird. Bei Wasserstoff-Luft-Gemischen verschiebt sich das Maximum dagegen
deutlich auf die fette Seite zu $\varphi \approx 1{,}7$
\cite{Law_2006, Egolfopoulos_1990, Dahoe_2005}. Ursache hierfür ist die hohe
Diffusivität des Wasserstoffmoleküls: In fetten Gemischen diffundiert
Wasserstoff schneller in die Reaktionszone hinein, als Wärme aus ihr
abgeführt wird, sodass die Umsatzrate über den Punkt maximaler
Flammentemperatur hinaus weiter ansteigt.

Die Abhängigkeit von Vorwärmtemperatur und Druck wird üblicherweise durch den
Potenzansatz

\begin{equation}
	\label{eq:sl_skalierung}
	s_\mathrm{L}(T_\mathrm{u}, p) = s_{\mathrm{L},0}
	\left(\frac{T_\mathrm{u}}{T_0}\right)^{\alpha}
	\left(\frac{p}{p_0}\right)^{\beta}
\end{equation}

beschrieben, wobei $s_{\mathrm{L},0}$ der Wert im Referenzzustand
($T_0 = 298\,$K, $p_0 = 1\,$bar) ist. Für Wasserstoff-Luft-Gemische liegt der
Temperaturexponent typischerweise bei $\alpha \approx 1{,}5\ldots2$, der
Druckexponent ist negativ ($\beta < 0$) \cite{Verhelst_2009}. Die
Druckabhängigkeit folgt unmittelbar aus der in
Abschnitt \ref{subsec:reaktionspfade} beschriebenen Konkurrenz zwischen
Kettenverzweigung und Kettenabbruch.

\section{Flammenstreckung und thermodiffusive Instabilität}\label{sec:streckung}

\subsection{Streckungsrate und Markstein-Zahl}\label{subsec:markstein}

Reale Flammen sind gekrümmt und einem Dehnungsfeld ausgesetzt. Die
Streckungsrate $\kappa$ beschreibt die relative zeitliche Änderung eines
Flammenflächenelements $A$,

\begin{equation}
	\label{eq:streckungsrate}
	\kappa = \frac{1}{A}\frac{\mathrm{d}A}{\mathrm{d}t}.
\end{equation}

Für schwache Streckung besteht ein linearer Zusammenhang zwischen der
gestreckten Flammengeschwindigkeit $s_\mathrm{L}^{\kappa}$ und der
ungestreckten Flammengeschwindigkeit $s_\mathrm{L}^{0}$,

\begin{equation}
	\label{eq:markstein}
	s_\mathrm{L}^{\kappa} = s_\mathrm{L}^{0} - \mathcal{L}\,\kappa
	\qquad\text{bzw.}\qquad
	\frac{s_\mathrm{L}^{\kappa}}{s_\mathrm{L}^{0}}
	= 1 - \mathrm{Ma}\cdot\mathrm{Ka},
\end{equation}

mit der Markstein-Länge $\mathcal{L}$, der Markstein-Zahl
$\mathrm{Ma} = \mathcal{L}/\delta_\mathrm{L}$ und der Karlovitz-Zahl
$\mathrm{Ka} = \kappa\,\delta_\mathrm{L}/s_\mathrm{L}^{0}$
\cite{Poinsot_2011, Matalon_2007}.

\subsection{Lewis-Zahl und zellulare Instabilität}\label{subsec:lewis}

Das Vorzeichen der Markstein-Zahl wird im Wesentlichen durch die Lewis-Zahl

\begin{equation}
	\label{eq:lewis}
	\mathrm{Le} = \frac{a}{D} = \frac{\lambda}{\rho\,c_p\,D}
\end{equation}

bestimmt, also durch das Verhältnis von thermischer zu molekularer
Diffusivität des Mangelreaktanden. Für $\mathrm{Le} > 1$ wirkt die Streckung
stabilisierend, für $\mathrm{Le} < 1$ destabilisierend.

Magere Wasserstoff-Luft-Flammen erreichen Lewis-Zahlen von etwa
$\mathrm{Le} \approx 0{,}3$, da Wasserstoff als Mangelreaktand wesentlich
schneller diffundiert, als Wärme geleitet wird. In einem konvex zum Frischgas
gekrümmten Flammenabschnitt fokussiert die Wasserstoffdiffusion den Brennstoff
lokal, während die Wärme divergent abgeführt wird. Die lokale
Flammentemperatur und damit die lokale Ausbreitungsgeschwindigkeit steigen an,
die Auslenkung verstärkt sich. Es entsteht eine selbstverstärkende, zellulare
Flammenstruktur \cite{Matalon_2007, Berger_2022}.

Diese thermodiffusive Instabilität überlagert sich der hydrodynamischen
Darrieus-Landau-Instabilität, die aus der Dichteänderung über die Flammenfront
resultiert und für alle Vormischflammen wirksam ist. Die praktische Konsequenz
ist erheblich: Durch die Vergrößerung der Flammenoberfläche kann die effektive
Ausbreitungsgeschwindigkeit einer zellularen Wasserstoffflamme die
eindimensional berechnete laminare Flammengeschwindigkeit um ein Vielfaches
übersteigen \cite{Frouzakis_2015, Berger_2022}. Eindimensionale Rechnungen
liefern in diesem Bereich also eine untere Schranke für die Rückschlagsneigung.
Dieser Umstand ist bei der Interpretation der Ergebnisse zwingend zu
berücksichtigen.
% TODO: Verweis auf das Ergebniskapitel ergaenzen, sobald dieses angelegt ist

\section{Turbulente Vormischverbrennung}\label{sec:turbulent}

Am untersuchten Brenner liegen turbulente Strömungsbedingungen vor. Die
turbulente Flammengeschwindigkeit $s_\mathrm{T}$ übersteigt $s_\mathrm{L}$,
weil die Turbulenz die Flammenfront faltet und dadurch deren Oberfläche
vergrößert. Nach Damköhler \cite{Damkoehler_1940} gilt im Bereich großskaliger
Turbulenz

\begin{equation}
	\label{eq:damkoehler}
	\frac{s_\mathrm{T}}{s_\mathrm{L}} = \frac{A_\mathrm{T}}{A_\mathrm{L}}
	\approx 1 + C\,\frac{u^{\prime}}{s_\mathrm{L}},
\end{equation}

mit der turbulenten Schwankungsgeschwindigkeit $u^{\prime}$ und einer
empirischen Konstanten $C$.

Zur Einordnung des Verbrennungsregimes dienen die Damköhler-Zahl und die
Karlovitz-Zahl,

\begin{equation}
	\label{eq:da_ka}
	\mathrm{Da} = \frac{\tau_\mathrm{t}}{\tau_\mathrm{c}}
	= \frac{l_\mathrm{t}/u^{\prime}}{\delta_\mathrm{L}/s_\mathrm{L}},
	\qquad
	\mathrm{Ka} = \frac{\tau_\mathrm{c}}{\tau_\eta}
	= \left(\frac{\delta_\mathrm{L}}{\eta}\right)^{2},
\end{equation}

worin $l_\mathrm{t}$ das integrale Längenmaß und $\eta$ das
Kolmogorov-Längenmaß bezeichnen. Im Borghi-Diagramm \cite{Borghi_1985, Peters_2000}
lassen sich damit die Regime der gefalteten Flammenfronten
($\mathrm{Ka} < 1$), der dünnen Reaktionszonen
($1 < \mathrm{Ka} < 100$) und der gerührten Reaktionszonen abgrenzen.

Die Bestimmung von $\delta_\mathrm{L}$ und $s_\mathrm{L}$ aus den in
Abschnitt \ref{sec:cantera} beschriebenen Rechnungen erlaubt es somit, die
Betriebspunkte des Prüfstands im Borghi-Diagramm einzuordnen.

\section{Mechanismen des Flammenrückschlags}\label{sec:flashback}

Ein Flammenrückschlag tritt ein, wenn die Flamme sich lokal schneller stromauf
bewegt, als das Frischgas sie stromab transportiert. Die allgemeine Bedingung
lautet

\begin{equation}
	\label{eq:flashback_bedingung}
	s_\mathrm{T} > u_\mathrm{lokal}.
\end{equation}

Entscheidend ist der Zusatz \emph{lokal}: Der Rückschlag wird nicht durch die
mittlere Strömungsgeschwindigkeit ausgelöst, sondern durch das Vorhandensein
von Bereichen mit reduzierter Geschwindigkeit. Je nachdem, wodurch ein solcher
Bereich entsteht, unterscheidet man die im Folgenden beschriebenen
Mechanismen \cite{Kalantari_2017, Lieuwen_2008}.

\subsection{Rückschlag in der wandnahen Grenzschicht}\label{subsec:blf}

In der wandnahen Grenzschicht sinkt die Strömungsgeschwindigkeit zur Wand hin
auf null ab. Gleichzeitig wird die Flamme durch den Wärmeverlust an die Wand
gelöscht, sobald sie sich der Wand auf weniger als den Löschabstand
$\delta_\mathrm{q}$ nähert. Aus dem Gleichgewicht beider Effekte folgt der von
Lewis und von Elbe eingeführte kritische Geschwindigkeitsgradient

\begin{equation}
	\label{eq:kritischer_gradient}
	g_\mathrm{c} = \left.\frac{\partial u}{\partial y}\right|_\mathrm{Wand,krit}
	\approx \frac{s_\mathrm{L}}{\delta_\mathrm{q}}.
\end{equation}

Ein Rückschlag wird demnach erwartet, sobald der tatsächliche Wandgradient
unter $g_\mathrm{c}$ fällt. Wegen des kleinen Löschabstands von Wasserstoff
(Tabelle \ref{tab:stoffwerte}) liegt $g_\mathrm{c}$ um mehr als eine
Größenordnung über dem Wert von Methan.

Neuere Untersuchungen zeigen allerdings, dass dieses Modell die
Rückschlagsneigung systematisch unterschätzt. Eichler und Sattelmayer
\cite{Eichler_2012} sowie Baumgartner et al.\ \cite{Baumgartner_2016} weisen
mittels hochauflösender Geschwindigkeitsmessungen nach, dass die
Wärmefreisetzung der annähernden Flamme über den Dichtesprung einen
positiven Druckgradienten stromauf erzeugt, der die Grenzschicht ablöst. Die
Flamme schafft sich damit selbst das Rückströmgebiet, in dem sie sich
fortbewegt; der Vorgang ist folglich rückgekoppelt und nicht durch das
ungestörte Strömungsfeld allein beschreibbar.

Endres und Sattelmayer \cite{Endres_2018} bestätigen dieses Bild numerisch und
leiten daraus ein handhabbares Kriterium ab: Der Rückschlag setzt ein, sobald
die Ausdehnung der lokalen Ablösegebiete den Löschabstand der Flamme deutlich
überschreitet. Damit tritt an die Stelle des rein kinematischen Vergleichs von
Gleichung \eqref{eq:kritischer_gradient} eine geometrische Bedingung, die
Wärmefreisetzung und Strömungsablösung miteinander verknüpft. Für die
numerische Vorhersage bedeutet dies, dass sowohl die differentielle Diffusion
als auch die Thermodiffusion im Speziestransport berücksichtigt werden
müssen \cite{Endres_2018}.

\subsection{Verbrennungsinduziertes Wirbelaufplatzen}\label{subsec:civb}

In drallstabilisierten Brennern kann die Wärmefreisetzung der Flamme das
Wirbelaufplatzen stromauf verlagern. Die Flamme reduziert über die
Volumenexpansion die axiale Impulsstromdichte im Wirbelkern und schwächt
gleichzeitig durch Baroklinität die azimutale Vortizität. Beides begünstigt
die Ausbildung eines Rückströmgebiets weiter stromauf, in das die Flamme
nachrückt; der Vorgang setzt sich selbstverstärkend fort. Kröner et al.\
\cite{Kroner_2003} und Fritz et al.\ \cite{Fritz_2004} haben diesen als CIVB
bezeichneten Mechanismus systematisch untersucht und quantitative
Rückschlagsgrenzen angegeben.

\subsection{Rückschlag in der Kernströmung}\label{subsec:core}

Übersteigt die turbulente Flammengeschwindigkeit die axiale
Strömungsgeschwindigkeit im Kern der Strömung, so propagiert die Flamme
unabhängig von Wand- und Dralleffekten stromauf. Dieser Mechanismus wird
insbesondere bei hohen Äquivalenzverhältnissen und geringen
Strömungsgeschwindigkeiten relevant. Da $s_\mathrm{T}$ nach
Gleichung \eqref{eq:damkoehler} unmittelbar mit $s_\mathrm{L}$ skaliert, ist
die laminare Flammengeschwindigkeit hier die maßgebliche Eingangsgröße --
weshalb die eindimensionalen Rechnungen dieser Arbeit für genau diesen
Mechanismus die belastbarste Aussage liefern. Ali und Bastiaans
\cite{Ali_2026} untersuchen den Vorgang für ultramagere Freistrahlflammen und
zeigen, dass die thermodiffusiv verstärkte Flammenausbreitung dort auch bei
sehr kleinen Äquivalenzverhältnissen zum Rückschlag führen kann. Fruzza et al.\
\cite{Fruzza_2025} weisen darauf hin, dass eine Reduktion auf zwei Dimensionen
die Flammendynamik verfälscht -- ein Hinweis, der für eine spätere
CFD-Betrachtung unmittelbar relevant ist.

\subsection{Akustisch und instabilitätsgetriebener Rückschlag}\label{subsec:acoustic}

Verbrennungsinstabilitäten koppeln Wärmefreisetzung und Druckschwankungen im
Brennraum. Die dabei auftretenden Geschwindigkeitsschwankungen können die
momentane Anströmgeschwindigkeit periodisch so weit absenken, dass die Flamme
kurzzeitig stromauf läuft. Dieser Mechanismus ist transient und lässt sich mit
stationären Betrachtungen grundsätzlich nicht erfassen.

\subsection{Kenngrößen zur Bewertung}\label{subsec:kenngroessen}

Zur dimensionslosen Beschreibung der Rückschlagsgrenze hat sich die
Flammen-Péclet-Zahl

\begin{equation}
	\label{eq:peclet}
	\mathrm{Pe} = \frac{g_\mathrm{c}\,\delta_\mathrm{L}}{s_\mathrm{L}}
\end{equation}

etabliert, die eine Übertragung zwischen unterschiedlichen Geometrien und
Betriebsbedingungen erlaubt \cite{Kalantari_2017}. Sowohl $g_\mathrm{c}$ als
auch $\delta_\mathrm{L}$ und $s_\mathrm{L}$ sind aus den in dieser Arbeit
durchgeführten Rechnungen beziehungsweise aus der Prüfstandsgeometrie
zugänglich.

\section{Numerische Modellierung mit Cantera}\label{sec:cantera}

\subsection{Modellgleichungen}\label{subsec:modellgleichungen}

Cantera \cite{Cantera_2023} ist eine quelloffene Programmbibliothek zur
Lösung von Problemen der chemischen Kinetik, der Thermodynamik und des
Stofftransports. Für die vorliegende Arbeit wird das Modell der frei
propagierenden, adiabaten eindimensionalen Vormischflamme
(\texttt{FreeFlame}) verwendet.

Im flammenfesten Bezugssystem lauten die stationären Erhaltungsgleichungen
\cite{Poinsot_2011, Cantera_2023}:

\begin{equation}
	\label{eq:kontinuitaet}
	\dot{m} = \rho\,u = \mathrm{const.},
\end{equation}

\begin{equation}
	\label{eq:spezies}
	\dot{m}\,\frac{\mathrm{d}Y_k}{\mathrm{d}x}
	+ \frac{\mathrm{d}}{\mathrm{d}x}\left(\rho\,Y_k\,V_k\right)
	= \dot{\omega}_k\,W_k,
	\qquad k = 1,\ldots,K,
\end{equation}

\begin{equation}
	\label{eq:energie}
	\dot{m}\,c_p\,\frac{\mathrm{d}T}{\mathrm{d}x}
	= \frac{\mathrm{d}}{\mathrm{d}x}\left(\lambda\,\frac{\mathrm{d}T}{\mathrm{d}x}\right)
	- \sum_{k=1}^{K}\rho\,Y_k\,V_k\,c_{p,k}\,\frac{\mathrm{d}T}{\mathrm{d}x}
	- \sum_{k=1}^{K} \dot{\omega}_k\,h_k\,W_k.
\end{equation}

Hierin sind $V_k$ die Diffusionsgeschwindigkeit, $W_k$ die molare Masse und
$h_k$ die spezifische Enthalpie der Spezies $k$. Der Eigenwert des Problems
ist der Massenstrom $\dot{m}$, aus dem die laminare Flammengeschwindigkeit
unmittelbar folgt:

\begin{equation}
	\label{eq:sl_aus_massenstrom}
	s_\mathrm{L} = \frac{\dot{m}}{\rho_\mathrm{u}}.
\end{equation}

\subsection{Transportmodell}\label{subsec:transport}

Cantera stellt für die Berechnung der Diffusionsgeschwindigkeiten mehrere
Modelle bereit. Das Modell \texttt{mixture-averaged} berechnet für jede Spezies
einen effektiven Diffusionskoeffizienten gegenüber der Mischung, das Modell
\texttt{multicomponent} löst dagegen die vollständige Matrix der binären
Diffusionskoeffizienten und erlaubt zusätzlich die Berücksichtigung der
Thermodiffusion (Soret-Effekt).

Gerade für Wasserstoff ist diese Unterscheidung nicht unerheblich: Wegen der
geringen molaren Masse von \ce{H2} und \ce{H} liefert das gemittelte Modell
insbesondere bei mageren Gemischen merklich abweichende Ergebnisse. In dieser
Arbeit wird daher der Einfluss des Transportmodells auf $s_\mathrm{L}$
explizit untersucht.

\subsection{Numerische Umsetzung und Gitterunabhängigkeit}\label{subsec:numerik}

Das Gleichungssystem wird auf einem adaptiven Gitter mit einem gedämpften
Newton-Verfahren gelöst. Die Gitterverfeinerung wird über die Parameter
\texttt{ratio}, \texttt{slope} und \texttt{curve} gesteuert, welche das
maximale Verhältnis benachbarter Zellweiten sowie die maximal zulässige
Änderung von Steigung und Krümmung der Lösungsvariablen zwischen zwei
Gitterpunkten begrenzen.

Da die berechnete Flammengeschwindigkeit von diesen Parametern abhängt, ist
vor der eigentlichen Parameterstudie eine Gitterunabhängigkeitsstudie
durchzuführen. Ebenso ist sicherzustellen, dass das Rechengebiet hinreichend
groß gewählt wird, damit die Randbedingungen die Lösung nicht beeinflussen.
Beide Punkte werden im Methodikkapitel behandelt.
% TODO: Verweis auf das Methodikkapitel ergaenzen, sobald dieses angelegt ist

% TODO: Verweis auf das Methodikkapitel korrigieren, sobald dieses angelegt ist
